\documentclass[a4paper,10pt]{report}\input{../0.Préambule/Préambule_cours_prof}\begin{document}

\pagestyle{empty}
\newpage
\noindent NOM, Prénom et classe : \dotfill

\section*{Les cercles}
\addcontentsline{toc}{section}{Les cercles}

\noindent \textit{Il sera tenu compte dans la notation de la propreté ainsi que de la justification apportée à chacune des réponses.}

\noindent \textit{Le barème est donné à titre indicatif. Il pourra être modifié ultérieurement.}

\noindent \textit{L'usage de la calculatrice est \textbf{interdite}.} \hfill \textit{Durée : 55 minutes}

\medskip
\textbf{Compétence 1 :} Je connais le vocabulaire autour du cercle.

\medskip
\textbf{Compétence 2 :} Je suis capable de me servir du compas.


\medskip
\paragraph{Exercice 1}: \hfill \textbf{5 points}

\noindent \begin{minipage}{6cm}
\begin{center}
\begin{tikzpicture}[line cap=round,line join=round,>=triangle 45,x=1.0cm,y=1.0cm,scale=0.7]
\clip(-3.8,-3.5) rectangle (5.,5.6);
\draw [line width=2.pt] (1.,1.) circle (4.cm);
\draw [line width=2.pt] (1.,1.)-- (2.8613178343399155,4.540550228363971);
\draw [line width=2.pt] (1.8011759373897895,2.7818577713576333) -- (2.013608951091628,2.670178701297238);
\draw [line width=2.pt] (1.8477088832482877,2.8703715270667325) -- (2.0601418969501264,2.758692457006337);
\draw [line width=2.pt] (1.,1.)-- (3.93689907756479,-1.7156258593921008);
\draw [line width=2.pt] (2.5132070760945977,-0.23576063412670578) -- (2.3502695245310723,-0.41197457878059257);
\draw [line width=2.pt] (2.5866295530337173,-0.30365128061150815) -- (2.4236920014701924,-0.4798652252653949);
\draw [line width=2.pt] (1.,1.)-- (-1.9368990775647903,3.7156258593921008);
\draw [line width=2.pt] (-0.5132070760945981,2.2357606341267053) -- (-0.35026952453107185,2.411974578780593);
\draw [line width=2.pt] (-0.586629553033718,2.3036512806115077) -- (-0.4236920014701917,2.4798652252653954);
\draw [line width=2.pt] (-1.9368990775647903,3.7156258593921008)-- (-3.,1.);
\draw [line width=2.pt] (-3.,1.)-- (2.8613178343399155,4.540550228363971);
\draw [line width=2.pt] (1.,1.)-- (2.088233710223408,-2.849122938012424);
\draw [line width=2.pt] (1.6459876218742842,-0.8438004209743543) -- (1.4150402455935376,-0.9090944435877589);
\draw [line width=2.pt] (1.6731934646298703,-0.9400284944246653) -- (1.4422460883491237,-1.00532251703807);
\begin{scriptsize}
\draw [color=blue] (1.,1.)-- ++(-3.5pt,0 pt) -- ++(7.0pt,0 pt) ++(-3.5pt,-3.5pt) -- ++(0 pt,7.0pt);
\draw[color=blue] (0.52,0.65) node {\large{$E$}};
\draw[color=black] (-0.84,5.35) node {\large{$\mathcal{C}$}};
\draw [color=blue] (2.8613178343399155,4.540550228363971)-- ++(-3.5pt,0 pt) -- ++(7.0pt,0 pt) ++(-3.5pt,-3.5pt) -- ++(0 pt,7.0pt);
\draw[color=blue] (3.,4.99) node {\large{$B$}};
\draw [color=blue] (3.93689907756479,-1.7156258593921008)-- ++(-3.5pt,0 pt) -- ++(7.0pt,0 pt) ++(-3.5pt,-3.5pt) -- ++(0 pt,7.0pt);
\draw[color=blue] (4.54,-1.79) node {\large{$D$}};
\draw [color=blue] (3.93689907756479,-1.7156258593921008)-- ++(-3.5pt,0 pt) -- ++(7.0pt,0 pt) ++(-3.5pt,-3.5pt) -- ++(0 pt,7.0pt);
\draw [color=blue] (-1.9368990775647903,3.7156258593921008)-- ++(-3.5pt,0 pt) -- ++(7.0pt,0 pt) ++(-3.5pt,-3.5pt) -- ++(0 pt,7.0pt);
\draw[color=blue] (-2.3,4.39) node {\large{$A$}};
\draw [color=blue] (-3.,1.)-- ++(-3.5pt,0 pt) -- ++(7.0pt,0 pt) ++(-3.5pt,-3.5pt) -- ++(0 pt,7.0pt);
\draw[color=blue] (-3.58,1.25) node {\large{$C$}};
\draw [color=blue] (2.088233710223408,-2.849122938012424)-- ++(-3.5pt,0 pt) -- ++(7.0pt,0 pt) ++(-3.5pt,-3.5pt) -- ++(0 pt,7.0pt);
\draw[color=blue] (2.62,-3.19) node {\large{$F$}};
\end{scriptsize}
\end{tikzpicture}
\end{center}
\end{minipage}
\begin{minipage}{11cm}
\begin{enumerate}
\item Complète les phrases suivants en utilisant les mots :

\begin{center}
\noindent "cercle" \quad "corde" \quad "rayon" \quad "centre" \quad "diamètre" \quad "milieu"
\end{center}

\begin{enumerate}
\item[•] Le \dots ~\dots ~\dots~\dots~\dots~\dots $(\mathcal{C})$ de \dots ~\dots ~\dots~\dots~\dots~\dots  $E$ passe par les points $A$, $B$, $C$, $D$ et $F$.\medskip
\item[•] Le segment $[EF]$ est un \dots ~\dots ~\dots~\dots~\dots~\dots de ce cercle.\medskip
\item[•] Le segment $[AC]$ est une \dots ~\dots ~\dots~\dots~\dots~\dots de ce cercle.\medskip
\item[•] $E$ est le \dots ~\dots ~\dots~\dots~\dots~\dots   du \dots ~\dots ~\dots~\dots~\dots~\dots $[AD]$.\medskip
\end{enumerate}
\end{enumerate}
\end{minipage}

\begin{enumerate}
\setcounter{enumi}{1}
\item Si la longueur AD=8 cm, quelle est la longueur du segment [AE] ?

\medskip\noindent\grandscarreaux{\linewidth}{1.6cm}
\end{enumerate}

\medskip
\paragraph{Exercice 2}: \hfill \textbf{4 points}

\medskip
\noindent Recopier la figure de gauche sur le quadrillage de droite

\noindent \begin{minipage}{8cm}
\begin{tikzpicture}[line cap=round,line join=round,>=triangle 45,x=1.0cm,y=1.0cm,scale=0.9]
\draw [color=lightgray,, xstep=1.0cm,ystep=1.0cm] (-1.5,-1.5) grid (7.5,7.5);
\clip(-1.5,-1.5) rectangle (7.5,7.5);
\draw [line width=1.2pt] (3.,3.) circle (4.cm);
\draw [shift={(3.,6.)},line width=1.2pt]  plot[domain=1.5707963267948963:4.71238898038469,variable=\t]({1.*1.*cos(\t r)+0.*1.*sin(\t r)},{0.*1.*cos(\t r)+1.*1.*sin(\t r)});
\draw [shift={(3.,0.)},line width=1.2pt]  plot[domain=-1.5707963267948966:1.5707963267948968,variable=\t]({1.*1.*cos(\t r)+0.*1.*sin(\t r)},{0.*1.*cos(\t r)+1.*1.*sin(\t r)});
\draw [shift={(6.,3.)},line width=1.2pt]  plot[domain=0.:3.141592653589793,variable=\t]({1.*1.*cos(\t r)+0.*1.*sin(\t r)},{0.*1.*cos(\t r)+1.*1.*sin(\t r)});
\draw [shift={(0.,3.)},line width=1.2pt]  plot[domain=3.141592653589793:6.283185307179586,variable=\t]({1.*1.*cos(\t r)+0.*1.*sin(\t r)},{0.*1.*cos(\t r)+1.*1.*sin(\t r)});
\draw [shift={(3.,5.)},line width=1.2pt]  plot[domain=-1.5707963267948966:1.5707963267948963,variable=\t]({1.*2.*cos(\t r)+0.*2.*sin(\t r)},{0.*2.*cos(\t r)+1.*2.*sin(\t r)});
\draw [shift={(5.,3.)},line width=1.2pt]  plot[domain=3.141592653589793:6.283185307179586,variable=\t]({1.*2.*cos(\t r)+0.*2.*sin(\t r)},{0.*2.*cos(\t r)+1.*2.*sin(\t r)});
\draw [shift={(3.,1.)},line width=1.2pt]  plot[domain=1.5707963267948968:4.71238898038469,variable=\t]({1.*2.*cos(\t r)+0.*2.*sin(\t r)},{0.*2.*cos(\t r)+1.*2.*sin(\t r)});
\draw [shift={(1.,3.)},line width=1.2pt]  plot[domain=0.:3.141592653589793,variable=\t]({1.*2.*cos(\t r)+0.*2.*sin(\t r)},{0.*2.*cos(\t r)+1.*2.*sin(\t r)});
\draw [shift={(2.,3.)},line width=1.2pt]  plot[domain=0.:3.141592653589793,variable=\t]({1.*1.*cos(\t r)+0.*1.*sin(\t r)},{0.*1.*cos(\t r)+1.*1.*sin(\t r)});
\draw [shift={(3.,4.)},line width=1.2pt]  plot[domain=-1.5707963267948966:1.5707963267948966,variable=\t]({1.*1.*cos(\t r)+0.*1.*sin(\t r)},{0.*1.*cos(\t r)+1.*1.*sin(\t r)});
\draw [shift={(4.,3.)},line width=1.2pt]  plot[domain=3.141592653589793:6.283185307179586,variable=\t]({1.*1.*cos(\t r)+0.*1.*sin(\t r)},{0.*1.*cos(\t r)+1.*1.*sin(\t r)});
\draw [shift={(3.,2.)},line width=1.2pt]  plot[domain=1.5707963267948966:4.71238898038469,variable=\t]({1.*1.*cos(\t r)+0.*1.*sin(\t r)},{0.*1.*cos(\t r)+1.*1.*sin(\t r)});
\end{tikzpicture}
\end{minipage}
\hspace{6mm}
\begin{minipage}{8cm}
\begin{tikzpicture}[line cap=round,line join=round,>=triangle 45,x=1.0cm,y=1.0cm,scale=0.9]
\draw [color=lightgray,, xstep=1.0cm,ystep=1.0cm] (-1.5,-1.5) grid (7.5,7.5);
\end{tikzpicture}
\end{minipage}

\medskip
\paragraph{Exercice 3}: \hfill \textbf{5 points}

\noindent Compléter les égalités suivantes :
\begin{enumerate}
\renewcommand{\arraystretch}{1.5}
\begin{tabularx}{\linewidth}{*{3}{>{\arraybackslash}X}}
\item $7,54 \times 0,001 = $ \dots ~\dots ~\dots ~\dots ~\dots&
\item $185 \times $ \dots ~\dots ~\dots ~\dots ~\dots $ = 1,85$ &
\item $0,72 \times 100 =$ \dots ~\dots ~\dots ~\dots ~\dots \\
\item $0,1052 \times $ \dots ~\dots ~\dots ~\dots ~\dots $ = 105,2$ &
\item $0,084 = $\dots ~\dots ~\dots ~\dots ~\dots $ \times 0,84$ &\\
\end{tabularx}
\end{enumerate}

\medskip
\paragraph{Exercice 4}: \hfill \textbf{4 points}

Sur la figure ci-dessus, effectue les tracés demandés.
\begin{enumerate}
\item Trace le cercle de centre $A$ et de rayon $3$cm.
\item Trace le cercle de rayon $[BC]$ et de centre $C$.
\item Trace un cercle de centre $D$ et de diamètre $8$cm.
\item Trace un cercle de rayon $FB$ et de centre $E$.
\end{enumerate}

\vspace{1cm}
\begin{center}
\begin{tikzpicture}[line cap=round,line join=round,>=triangle 45,x=1.0cm,y=1.0cm]
\clip(0.5,1.5) rectangle (12.5,8.6);
\draw [line width=1.2pt] (9.,8.)-- (12.,5.);
\draw [line width=1.2pt] (3.,8.)-- (7.,3.);
\begin{scriptsize}
\draw [color=blue] (1.,6.)-- ++(-3.5pt,0 pt) -- ++(7.0pt,0 pt) ++(-3.5pt,-3.5pt) -- ++(0 pt,7.0pt);
\draw[color=blue] (1.14,6.45) node {\large{$A$}};
\draw [color=blue] (9.,8.)-- ++(-3.5pt,0 pt) -- ++(7.0pt,0 pt) ++(-3.5pt,-3.5pt) -- ++(0 pt,7.0pt);
\draw[color=blue] (9.14,8.45) node {\large{$B$}};
\draw [color=blue] (12.,5.)-- ++(-3.5pt,0 pt) -- ++(7.0pt,0 pt) ++(-3.5pt,-3.5pt) -- ++(0 pt,7.0pt);
\draw[color=blue] (12.14,5.45) node {\large{$C$}};
\draw [color=blue] (4.,2.)-- ++(-3.5pt,0 pt) -- ++(7.0pt,0 pt) ++(-3.5pt,-3.5pt) -- ++(0 pt,7.0pt);
\draw[color=blue] (4.14,2.45) node {\large{$D$}};
\draw [color=blue] (3.,8.)-- ++(-3.5pt,0 pt) -- ++(7.0pt,0 pt) ++(-3.5pt,-3.5pt) -- ++(0 pt,7.0pt);
\draw[color=blue] (3.14,8.45) node {\large{$E$}};
\draw [color=blue] (7.,3.)-- ++(-3.5pt,0 pt) -- ++(7.0pt,0 pt) ++(-3.5pt,-3.5pt) -- ++(0 pt,7.0pt);
\draw[color=blue] (7.14,3.45) node {\large{$F$}};
\end{scriptsize}
\end{tikzpicture}
\end{center}

\medskip
\paragraph{Exercice 5}: \hfill \textbf{2 points}

Faire une rosace.

\end{document}